\documentclass[11pt]{article}
\usepackage[T1]{fontenc}
\usepackage{amsmath}
\usepackage{amssymb}
\usepackage{amsthm}
\usepackage[hidelinks]{hyperref}
\newtheorem{theorem}{Theorem}[section]
\newtheorem{lemma}[theorem]{Lemma}
\newtheorem{corollary}[theorem]{Corollary}
\theoremstyle{definition}
\newtheorem{definition}[theorem]{Definition}
\begin{document}
\begin{center}
{\LARGE Natural number addition : consequences}
\end{center}
\section{Commutativity : consequences}
\label{ll:main:consequences}
\begin{theorem}
\label{ll:main:comm-under-hypothesis}
For every natural number \(x\) and natural number \(y\), if \(x = y\), then \(x + y = y + x\).
\end{theorem}
\begin{proof}
Assume \(h\).
The goal follows from \(\operatorname{add_comm}(x, y)\).
\end{proof}
\begin{theorem}
\label{ll:main:comm-rewrite}
For every natural number \(x\) and natural number \(y\), if \(x = y\), then \(x + x = y + x\).
\end{theorem}
\begin{proof}
Assume \(h\).
Rewrite the goal using \(\rightarrow h\).
\end{proof}
\begin{theorem}
\label{ll:main:double-divides}
For every natural number \(n\), \(\operatorname{divides}(2, \operatorname{double}(n))\).
\end{theorem}
\begin{proof}
Use \(n\) as the witness.
The goal follows from \(\operatorname{eqsymm}(\operatorname{twomul}(n))\).
\end{proof}
\begin{theorem}
\label{ll:main:even-divides}
For every natural number \(n\), if \(\operatorname{even}(n)\), then \(\operatorname{divides}(2, n)\).
\end{theorem}
\begin{proof}
Assume \(h\).
Consider the cases of \(h\).
Case \(existsintro\) with \(k\), \(hk\):
Use \(k\) as the witness.
Rewrite the goal using \(\rightarrow hk\).
The goal follows from \(\operatorname{eqsymm}(\operatorname{twomul}(k))\).
\end{proof}
\section{Excluded middle}
\label{ll:main:excluded-middle}
\begin{theorem}
\label{ll:main:even-or-not}
For every natural number \(n\), \(\operatorname{even}(n)\) or not \(\operatorname{even}(n)\).
\end{theorem}
\begin{proof}
The goal follows from \(\operatorname{em}(\operatorname{even}(n))\).
\end{proof}
\begin{theorem}
\label{ll:main:positive-or-zero}
For every natural number \(n\), \(\operatorname{positive}(n)\) or \(n = 0\).
\end{theorem}
\begin{proof}
Consider the cases of \(n\).
Case zero:
Select the right alternative.
The goal follows by reflexivity.
Case \(succ\) with \(m\):
Select the left alternative.
The goal follows from \(\operatorname{successor_positive}(m)\).
\end{proof}
\end{document}
